% GENERATED FILE. Edit canonical lessons, then run npm run generate:books.
% canonical-source-hash: fnv1a64:82bdb0d58ec07d58
% canonical-lessons: IT-C485-slogarsi, IT-C485-crampo, IT-C485-gonfio, IT-C485-garza, IT-C485-fasciare

\chapter{Una brutta storta}
\label{ch:it-una-brutta-storta}
\hlchaptermodality{485}

\begin{chapteropening}
\textbf{By the end of this chapter:} \emph{I can describe a sprain, a cramp and how it is bandaged.}

Complete the last lesson of chapter 485: I can describe a sprain, a cramp and how it is bandaged.
\end{chapteropening}

\section[slogarsi]{slogarsi — to sprain, to dislocate}

\label{lesson:IT-C485-slogarsi}



\begin{quote}
Before the new one: say the Italian for to put in a plaster cast, then the Italian for to limp.
\end{quote}

\subsection*{You'll want to know: slogarsi}
\textbf{slogarsi} — \textquotedblleft{}to sprain, to dislocate\textquotedblright{}.

\begin{quote}
Giocando a tennis mi sono slogato la caviglia.
\end{quote}

\begin{quote}
\emph{Playing tennis I sprained my ankle.}
\end{quote}

\emph{Slogarsi la caviglia} — to sprain one's ankle; in speech also \emph{prendere una storta}.

\begin{cousinweb}
\emph{slogarsi} is built from \emph{s-} \textquotedblleft{}out of\textquotedblright{} and \emph{luogo} \textquotedblleft{}place\textquotedblright{}, inherited from Latin \emph{locus} (Lewis \& Short): a joint put out of its place.
\end{cousinweb}

\begin{grammarlens}[title={What you've built}]
One more word for telling, explaining and describing at B1.
\end{grammarlens}

\subsection*{Your turn}
\begin{itemize}
\raggedright
  \item \emph{Say it:} \emph{slogarsi}
  \item \emph{Say it:} \emph{slogarsi}, once more
  \item \emph{Say it:} a sentence of your own with \emph{slogarsi}
  \item \emph{Recall:} say the Italian for to put in a plaster cast, then the Italian for to limp, then say \emph{slogarsi} again
\end{itemize}

\begin{tcolorbox}[breakable,colback=teal!4,colframe=teal!35!black,title={Before you move on}]
What does slogarsi mean? (\textquotedblleft{}To sprain, to dislocate\textquotedblright{}.) Say it once more.
\end{tcolorbox}
\section[il crampo]{il crampo — cramp}

\label{lesson:IT-C485-crampo}



\begin{quote}
Before the new one: say the Italian for to limp, then the Italian for to sprain.
\end{quote}

\subsection*{You'll want to know: il crampo}
\textbf{il crampo} — \textquotedblleft{}cramp\textquotedblright{}.

\begin{quote}
Di notte mi viene un crampo al polpaccio.
\end{quote}

\begin{quote}
\emph{At night I get a cramp in my calf.}
\end{quote}

\emph{Avere un crampo} — to have a cramp.

\begin{cousinweb}
\emph{crampo} is usually traced to a Germanic word, Proto-Germanic \emph{krampō} \textquotedblleft{}cramp, clasp\textquotedblright{} (EtymDB, under French \emph{crampe}).
\end{cousinweb}

\begin{grammarlens}[title={What you've built}]
One more word for telling, explaining and describing at B1.
\end{grammarlens}

\subsection*{Your turn}
\begin{itemize}
\raggedright
  \item \emph{Say it:} \emph{il crampo}
  \item \emph{Say it:} \emph{il crampo}, once more
  \item \emph{Say it:} a sentence of your own with \emph{il crampo}
  \item \emph{Recall:} say the Italian for to limp, then the Italian for to sprain, then say \emph{il crampo} again
\end{itemize}

\begin{tcolorbox}[breakable,colback=teal!4,colframe=teal!35!black,title={Before you move on}]
What does il crampo mean? (\textquotedblleft{}Cramp\textquotedblright{}.) Say it once more.
\end{tcolorbox}
\section[gonfio]{gonfio — swollen}

\label{lesson:IT-C485-gonfio}



\begin{quote}
Before the new one: say the Italian for to sprain, then the Italian for cramp.
\end{quote}

\subsection*{You'll want to know: gonfio}
\textbf{gonfio} — \textquotedblleft{}swollen\textquotedblright{}.

\begin{quote}
Ho la caviglia gonfia.
\end{quote}

\begin{quote}
\emph{My ankle is swollen.}
\end{quote}

The feminine is \emph{gonfia}. The verb is \emph{gonfiarsi} — to swell.

\begin{cousinweb}
\emph{gonfio} is built on \emph{gonfiare} \textquotedblleft{}to blow up, to swell\textquotedblright{}, inherited from Latin \emph{cōnflāre} \textquotedblleft{}to blow together\textquotedblright{} (Lewis \& Short), from \emph{flāre} \textquotedblleft{}to blow\textquotedblright{}.
\end{cousinweb}

\begin{grammarlens}[title={What you've built}]
One more word for telling, explaining and describing at B1.
\end{grammarlens}

\subsection*{Your turn}
\begin{itemize}
\raggedright
  \item \emph{Say it:} \emph{gonfio}
  \item \emph{Say it:} \emph{gonfio}, once more
  \item \emph{Say it:} a sentence of your own with \emph{gonfio}
  \item \emph{Recall:} say the Italian for to sprain, then the Italian for cramp, then say \emph{gonfio} again
\end{itemize}

\begin{tcolorbox}[breakable,colback=teal!4,colframe=teal!35!black,title={Before you move on}]
What does gonfio mean? (\textquotedblleft{}Swollen\textquotedblright{}.) Say it once more.
\end{tcolorbox}
\section[la garza]{la garza — gauze}

\label{lesson:IT-C485-garza}



\begin{quote}
Before the new one: say the Italian for cramp, then the Italian for swollen.
\end{quote}

\subsection*{You'll want to know: la garza}
\textbf{la garza} — \textquotedblleft{}gauze\textquotedblright{}.

\begin{quote}
Ho messo una garza pulita sulla ferita.
\end{quote}

\begin{quote}
\emph{I put a clean piece of gauze on the wound.}
\end{quote}

A thin cloth for wounds; the roll that holds it is \emph{la benda}.

\begin{cousinweb}
\emph{garza} comes from French \emph{gaze} \textquotedblleft{}gauze\textquotedblright{} (EtymDB), traditionally linked to a town on the eastern Mediterranean where the cloth was made.
\end{cousinweb}

\begin{grammarlens}[title={What you've built}]
One more word for telling, explaining and describing at B1.
\end{grammarlens}

\subsection*{Your turn}
\begin{itemize}
\raggedright
  \item \emph{Say it:} \emph{la garza}
  \item \emph{Say it:} \emph{la garza}, once more
  \item \emph{Say it:} a sentence of your own with \emph{la garza}
  \item \emph{Recall:} say the Italian for cramp, then the Italian for swollen, then say \emph{la garza} again
\end{itemize}

\begin{tcolorbox}[breakable,colback=teal!4,colframe=teal!35!black,title={Before you move on}]
What does la garza mean? (\textquotedblleft{}Gauze\textquotedblright{}.) Say it once more.
\end{tcolorbox}
\section[fasciare]{fasciare — to bandage, to wrap}

\label{lesson:IT-C485-fasciare}



\begin{quote}
Before the new one: say the Italian for swollen, then the Italian for gauze.
\end{quote}

\subsection*{You'll want to know: fasciare}
\textbf{fasciare} — \textquotedblleft{}to bandage, to wrap\textquotedblright{}.

\begin{quote}
Il medico mi ha fasciato la mano.
\end{quote}

\begin{quote}
\emph{The doctor bandaged my hand.}
\end{quote}

The bandage is \emph{la fascia} or \emph{la benda}.

\begin{cousinweb}
\emph{fasciare} is built on \emph{fascia} \textquotedblleft{}band\textquotedblright{}, inherited from Latin \emph{fascia} \textquotedblleft{}a band, a bandage\textquotedblright{} (Lewis \& Short).
\end{cousinweb}

\begin{grammarlens}[title={What you've built}]
One more word for telling, explaining and describing at B1.
\end{grammarlens}

\subsection*{Your turn}
\begin{itemize}
\raggedright
  \item \emph{Say it:} \emph{fasciare}
  \item \emph{Say it:} \emph{fasciare}, once more
  \item \emph{Say it:} a sentence of your own with \emph{fasciare}
  \item \emph{Recall:} say the Italian for swollen, then the Italian for gauze, then say \emph{fasciare} again
\end{itemize}

\begin{tcolorbox}[breakable,colback=teal!4,colframe=teal!35!black,title={Before you move on}]
What does fasciare mean? (\textquotedblleft{}To bandage, to wrap\textquotedblright{}.) Say it once more.
\end{tcolorbox}
